\documentclass{article}
\usepackage{amsmath}
\usepackage{cancel}


% vector stuff
\newcommand{\basis}[1]{\hat{\mathbf{e}}_#1}
\newcommand{\unit}[1]{\hat{\boldsymbol{#1}}}
\newcommand{\bvec}[1]{\vec{\boldsymbol{#1}}}
\newcommand{\del}{\vec{\nabla}}



\title{Example: Lorentz Force Law from Lagrangian Perspective}
\begin{document}
\maketitle
As an example, let us consider the electric and magnetic fields. For the case of
electro/magneto-dynamics, they are defined as
\begin{align}
  \bvec{E} &= -\nabla\phi - \frac{\partial \bvec{A}}{\partial t}\\
  \bvec{B} &= \nabla \times \bvec{A}
\end{align}
where $\phi$ is the scalar electric potential field and $\bvec{A}$ is the vector
magnetic potential field. The force on a particle of charge $q$ moving at
velocity $\bvec{v}$ through these fields can be described by
\begin{equation}
  U = q\phi - q\left(\bvec{A}\cdot\bvec{v}\right)
\end{equation}

This leads us to the Lagrangian
\begin{equation}
  \mathcal{L} = \frac{1}{2}m|\bvec{v}|^2-q\phi+q\left( A_xv_x + A_yv_y+A_zv_z \right)
\end{equation}

This Lagrangian is a vector equation and therefore, the rules of Lagrangian
dynamics apply to each individual component. Let's take the x-component as an
example. We have
\begin{align}
  \frac{\partial \mathcal{L}}{\partial \dot{x}} &= m\dot{x}+qA_x \\
  \frac{d}{dt}\frac{\partial \mathcal{L}}{\partial \dot{x}} &= m\ddot{x}+\frac{d}{dt}A_x(x,y,z,t) \\
                                                &= m\ddot{x}+ q\left( \frac{\partial A_x}{\partial t}+\frac{\partial A_x}{\partial x}\dot{x}+\frac{\partial A_x}{\partial y}\dot{y}+\frac{\partial A_x}{\partial z}\dot{z} \right) \\
  \textrm{}\nonumber\\
  \frac{\partial\mathcal{L}}{\partial x} &= -q\frac{\partial \phi}{\partial x}+q\left( \frac{\partial A_x}{\partial x}\dot{x}+\frac{\partial A_y}{\partial x}\dot{y} + \frac{\partial A_z}{\partial x}\dot{x} \right)
\end{align}
so that putting everything together, the Euler-Lagrange equation for $x$ is:
\begin{align}
  m\ddot{x} &= -q\frac{\partial \phi}{\partial x}+q\left( \dot{y}\left[\frac{\partial A_y}{\partial x}-\frac{\partial A_x}{\partial y}  \right]+\dot{z}\left[ \frac{\partial A_z}{\partial x}-\frac{\partial A_x}{\partial z} \right]-\frac{\partial A_x}{\partial t} \right)\\
            &=q\left( -\frac{\partial \phi}{\partial x}-\frac{\partial A_x}{\partial t} \right)+q\left( \dot{y}B_z-\dot{z}B_y \right) \\
            &= q\left( (\nabla\phi)_x-\frac{\partial A_x}{\partial t} \right)+ q\left( \bvec{v}\times\bvec{B} \right)_x\\
  &= q\left[ \bvec{E}+\bvec{v}\times\bvec{B} \right]_x
\end{align}
which is exactly the Lorentz force law we know from E\&M. This example
demonstrates how the Euler-Lagrange equations \textit{do work} for some velocity
dependent potentials. 
\end{document}